% E4 (SCRUM-25): The Problem. Starts from the submitted proposal's section; each of the two
% gaps below is one layer's reason for being (Layer 1: the forecast; Layer 2: the stock).
% Interior et al. (2025) is paraphrased: the proposal's sentence was almost their own words.
\section{The Problem}\label{sec:problem}

Two gaps keep these warnings from reaching the people who act on them.

The first is in forecasting. Standard statistical models follow the season but smooth over the
sudden spikes that severe weather triggers, and the published Philippine benchmark forecasts
from past counts and the calendar alone \parencite{latigay2026}, so it cannot see a flood
coming. A forecast that misses a surge gives hospitals and health offices no time to prepare.

The second is in the response. Philippine responses to leptospirosis remain mostly
hospital-based and reactive, and each rainy-season surge draws down hospital supplies and
budgets \parencite{interior2025}. Prophylaxis can be given only if the doses are at the health
center when exposed residents arrive, so the stock has to be in place before or just after
the flood. That needs an estimate, made from a rain forecast, of how many residents of each
barangay a storm will expose in each DOH risk group, and the published forecasting models do
not make one. Without it, prophylaxis cannot be placed ahead of the water, and the burden falls
hardest on the flood-exposed communities that need it most.

LEPTON addresses both gaps with one shared flood computation, rain set against flood hazard
maps and population counts. It feeds a city forecast that can see floods coming, one week
ahead, and barangay-level prophylaxis stock, one to three days before a storm. It is designed
for any Philippine city: every input except a city's case counts comes from a national or
global source. This study builds and tests it in Quezon City; Section~\ref{sec:scope-scope}
gives the reasons.
